\section{Open Source}

By making Zoo open-source, developers can build on top of Zoo and use the \$ZOO token as their projects' virtual utility token or build other games and NFT platforms parallel to the codebase.

This allows for the utility of the \$ZOO token to go beyond that which the project can support itself and become the virtual utility token of a potentially endless number of additional projects, creating enormous potential that is not possible without alignment to open-source ideology. Furthermore, the project's strong foundational support in the field of Education is another element that will make it shine through the use of open-source elements. This content will all be created and curated by the Zoo Foundation's resident Director of Education.

Zoo will be working with partners in conservation to establish content to be shared through the Zoo Foundation in an open-source format. This educational content, aligned to nationally recognized learning standards for grade levels K-12 and supplementary content, will be completely free to schools around the US. Schools will be encouraged to share the results of their time spent working with the materials in a repository of open-source materials found on the Zoo Foundation's website.

The development team is dedicated to improving the future of the world. The possibilities are endless and begin with community engagement.
